% !PRINTED 562
\weberpage{562}
If, besides the space $R_\xi$, we consider also a second space $R_x$,
in which the variables $x$ determined by (3) or (6) are the
coordinates, then to every point of the one space there corresponds
one and only one point of the other space. To a finite % [sic: the page prints ``endliche'' for ``endlichen''; the inflection error has no English counterpart]
region $G_\xi$ there corresponds a finite region $G_x$.

In order to ascertain the relation between the volumes $V_\xi$ and $V_x$,
we must treat the integral (8) according to the rules for the
transformation of multiple integrals, for which the general formula holds:

\weq{11}{\int\!\int\!\ldots\!\int d\xi_1\, d\xi_2 \ldots d\xi_n
= \int\!\int\!\cdots\!\int
\sqrt{\left(\Sigma\pm\frac{d\xi_1}{dx_1}\frac{d\xi_2}{dx_2}\cdots\frac{d\xi_n}{dx_n}\right)^{2}}\,
dx_1\, dx_2 \cdots dx_n\ \footnotemark.}

\footnotetext{The general formulation of this transformation is due to
\gesperrt{Jacobi} (De determinantibus functionalibus; \gesperrt{Crelle},
Bd. 22, 1891, % [sic: Weber prints 1891.]
gesammelte Werke, Bd. 3). The derivation is to be found
in most of the more detailed textbooks of the integral calculus, e.g.
\gesperrt{Serret-Harnack}, Bd. 2; \gesperrt{Lipschitz}, Lehrbuch der
Analysis Bd. 2; also \gesperrt{Baltzer}, Determinanten.}

If in this we insert for the functional determinant the value (4) or
(5), there results

\weq{12}{V_\xi = \sqrt{D}\ V_x,}

in which now the integral

\weq{13}{V_x = \int\!\int\!\cdots\!\int d x_1\, d x_2 \cdots d x_n}

refers to the region $G_x$.

The integral $V_x$ we may conceive as the limiting value of a sum of
elements $\Delta x_1 \Delta x_2 \ldots \Delta x_n$, and
if we give to all the increments $\Delta x_i$ the quantity $1 : t^{\frac1n}$,
then each of these elements will have the value $1 : t$.
If $Z_t$ denotes the number of these elements, then $V_x$ is the
limiting value which the ratio $Z_t : t$ approaches as $t$
increases indefinitely; thus

\weq{14}{V_x = \mathrm{Lim}_{t\,=\,\infty} \frac{Z_t}{t}.}

Now the number of elements $Z_t$ is just as great as the number of
points in the interior of $G_x$ whose coordinates have the
expression

\[
x_1 = h_1 t^{-\frac1n},\ x_2 = h_2 t^{-\frac1n},\ \ldots,\ x_n = h_n t^{-\frac1n}
\]

with integer $h_1, h_2, \ldots, h_n$, or again equal to the
%CONT>>
